\documentclass[11pt]{article}

\usepackage{amsmath,amssymb,amsfonts}
\usepackage{graphicx}
\usepackage{hyperref}
\usepackage{geometry}
\geometry{margin=2.5cm}

\title{Prime Quadruple Resonators and Logarithmic Spectral Fluctuations}
\author{Thomas Hoffbauer}
\date{\today}

\begin{document}
	
	\maketitle
	
	\begin{abstract}
		
		We introduce a family of observables associated with prime quadruples 
		$(p,p+2,p+6,p+8)$ arranged in mod-$12$ families $(E,A,B,C)$.  
		For each quadruple we define a resonance defect based on the geometric mean of three primes.  
		
		After removing the asymptotic expansion of this defect we analyze the residual fluctuations in the logarithmic variable $u=\log x$.  
		Numerical evidence suggests that these fluctuations exhibit oscillatory components compatible with frequencies associated with the imaginary parts of the nontrivial zeros of the Riemann zeta function.
		
		This motivates a heuristic spectral interpretation of prime quadruples as local arithmetic resonators and suggests a discrete operator model in the spirit of the Hilbert–Pólya program.
		
	\end{abstract}
	
	\section{Introduction}
	
	The connection between the distribution of prime numbers and the nontrivial zeros of the Riemann zeta function
	
	\[
	\zeta(s)
	\]
	
	is encoded in the explicit formulas of analytic number theory.  
	These formulas express fluctuations in prime counting functions as oscillatory contributions involving the imaginary parts $\gamma$ of the zeros
	
	\[
	\rho=\frac12+i\gamma .
	\]
	
	In this work we investigate whether similar oscillatory signatures appear in observables derived from prime quadruples
	
	\[
	(p,p+2,p+6,p+8).
	\]
	
	These quadruples form two mod-$12$ families
	
	\[
	(A,B,C,E),\qquad (C,E,A,B)
	\]
	
	which will play the role of two chiral sectors of the system.
	
	\section{Prime Quadruple Defects}
	
	Consider a prime quadruple
	
	\[
	(p,p+2,p+6,p+8).
	\]
	
	Define
	
	\[
	(a,b,c,e)=(p,p+2,p+6,p+8).
	\]
	
	We introduce the resonance defect
	
	\[
	\Delta_{ABCE}(a)
	=
	e-(abc)^{1/3}-\frac{16}{3}.
	\]
	
	For the second ordering
	
	\[
	(c,e,a,b)=(p,p+2,p+6,p+8)
	\]
	
	we define
	
	\[
	\Delta_{CEAB}(c)
	=
	e-(abc)^{1/3}+4 .
	\]
	
	These quantities measure the deviation between the largest prime of the quadruple and the geometric mean of the preceding three primes.
	
	\section{Asymptotic Expansion}
	
	For large primes the defect admits an expansion
	
	\[
	\Delta(x)
	=
	\frac{K_1}{x}
	+
	\frac{K_2}{x^2}
	+
	\frac{K_3}{x^3}
	+
	O(x^{-4}).
	\]
	
	Explicit calculation yields
	
	\[
	\Delta_{ABCE}(a)
	=
	\frac{28}{9a}
	-
	\frac{832}{81a^2}
	+
	\frac{10288}{243a^3}
	+
	O(a^{-4})
	\]
	
	and
	
	\[
	\Delta_{CEAB}(c)
	=
	\frac{52}{9c}
	-
	\frac{1624}{81c^2}
	+
	\frac{23008}{243c^3}
	+
	O(c^{-4}).
	\]
	
	Removing these smooth contributions isolates the fluctuating component.
	
	\section{Logarithmic Spectral Analysis}
	
	Define
	
	\[
	u=\log x
	\]
	
	and the detrended signal
	
	\[
	y(u)
	=
	\Delta(x)-\Delta_{\text{asym}}(x).
	\]
	
	If the fluctuations are influenced by the zeta spectrum we expect an expansion
	
	\[
	y(u)
	\approx
	\sum_{\gamma}
	A_\gamma
	\cos(\gamma u+\phi_\gamma).
	\]
	
	We therefore project the signal onto the basis
	
	\[
	\cos(\gamma \log x),\qquad
	\sin(\gamma \log x)
	\]
	
	using the first several hundred imaginary parts $\gamma$ of the zeta zeros.
	
	\section{Numerical Results}
	
	The projections produce amplitudes
	
	\[
	\mathcal A_\gamma
	=
	\sqrt{A_\gamma^2+B_\gamma^2}.
	\]
	
	Empirically we observe that certain frequencies appear repeatedly in both families of prime quadruples.
	
	Figures illustrate
	
	\begin{itemize}
		\item the detrended defect as a function of $\log x$
		\item the reconstruction using the strongest spectral modes
		\item the amplitude spectrum across $\gamma$
	\end{itemize}
	
	\section{Operator Heuristic}
	
	The oscillatory structure suggests interpreting prime quadruples as elements of a discrete state space
	
	\[
	\mathcal H = \text{span}\{|q\rangle\}
	\]
	
	where each $|q\rangle$ corresponds to a quadruple.
	
	Define a diagonal observable
	
	\[
	H_0|q\rangle=\Delta(q)|q\rangle .
	\]
	
	Couplings between nearby quadruples define an interaction operator $T$.
	
	The total operator
	
	\[
	H=H_0+T
	\]
	
	or in chiral form
	
	\[
	D=
	\begin{pmatrix}
		0 & T^\ast\\
		T & 0
	\end{pmatrix}
	\]
	
	may produce spectral statistics related to the imaginary parts of the zeta zeros.
	
	\section{Geometric Embedding}
	
	To visualize the internal structure we embed quadruples into an eight-dimensional space
	
	\[
	(\log p_1,\log p_2,\log p_3,\log p_4,\phi_E,\phi_A,\phi_B,\phi_C)
	\]
	
	where the phases encode the mod-$12$ prime families.
	
	Projection of this space reveals preferred directions and possible chiral separation between the two quadruple types.
	
	\section{Discussion}
	
	The proposed observables define a new arithmetic object associated with prime quadruples.
	
	While the observed oscillations do not constitute a proof of any spectral interpretation, they motivate further investigation of discrete operator models connecting prime configurations with zeta spectra.
	
	Future work includes
	
	\begin{itemize}
		\item extending the numerical range of quadruples
		\item studying stability of spectral peaks
		\item constructing explicit discrete operators on quadruple graphs
	\end{itemize}
	
	\section{Conclusion}
	
	Prime quadruples equipped with geometric-mean defects provide a new framework for studying oscillatory phenomena in prime distributions.
	
	The logarithmic spectral analysis suggests a possible resonance structure compatible with frequencies related to the imaginary parts of the Riemann zeta zeros.
	
	\begin{thebibliography}{9}
		
		\bibitem{Edwards}
		H. M. Edwards,
		\textit{Riemann's Zeta Function},
		Dover.
		
		\bibitem{Connes}
		A. Connes,
		Trace formula in noncommutative geometry.
		
		\bibitem{Odlyzko}
		A. Odlyzko,
		Tables of zeros of the Riemann zeta function.
		
	\end{thebibliography}
	
\end{document}